\documentclass[aspectratio=169]{beamer}
\usetheme{Madrid}
\usecolortheme{default}

% Packages
\usepackage{graphicx}
\usepackage{booktabs}
\usepackage{tikz}
\usepackage{multirow}
\usepackage{array}
\usepackage{xcolor}

% Theme customization
\setbeamertemplate{navigation symbols}{}
\setbeamertemplate{footline}[frame number]

% Title information
\title[Retinal Image Enhancement]{Deep Learning-Based Retinal Image Enhancement:\\Zeiss Visuscout to Clarus Quality}
\subtitle{Final Year Project Presentation}
\author{Your Name}
\institute{Your University\\Department of Computer Science/Engineering}
\date{\today}

\begin{document}

% Title slide
\begin{frame}
\titlepage
\end{frame}

% Outline
\begin{frame}{Outline}
\tableofcontents
\end{frame}

% Section 1: Introduction
\section{Introduction}

\begin{frame}{Clinical Context: Diabetic Retinopathy}
\begin{columns}[T]
\column{0.5\textwidth}
\textbf{Global Health Challenge:}
\begin{itemize}
    \item Diabetic retinopathy is a leading cause of blindness
    \item Early detection is critical for prevention
    \item Requires high-quality retinal imaging for accurate diagnosis
    \item \textcolor{red}{\textbf{Problem:}} Many clinics use lower-quality imaging systems
\end{itemize}

\column{0.5\textwidth}
\textbf{Diagnostic Requirements:}
\begin{itemize}
    \item Clear visualization of blood vessels
    \item Detection of microaneurysms
    \item Identification of hemorrhages
    \item Assessment of retinal abnormalities
\end{itemize}
\end{columns}

\vspace{0.5cm}
\centering
\textcolor{blue}{\large \textbf{Solution: AI-powered image enhancement to bridge the quality gap}}
\end{frame}

% Section 2: Problem Statement
\section{Problem Statement}

\begin{frame}{Camera Systems Comparison}
\begin{table}
\centering
\renewcommand{\arraystretch}{1.5}
\begin{tabular}{|l|c|c|}
\hline
\textbf{Specification} & \textbf{\textcolor{green!50!black}{Clarus (Reference)}} & \textbf{\textcolor{red!70!black}{Zeiss Visuscout (Input)}} \\
\hline\hline
\textbf{Field of View (FOV)} & \textcolor{green!50!black}{\huge \textbf{133°}} & \textcolor{red!70!black}{\huge \textbf{40°}} \\
\hline
\textbf{Resolution/DPI} & \textcolor{green!50!black}{\textbf{300 DPI}} & Lower resolution \\
\hline
\textbf{Coverage} & Wide retinal area & \textcolor{red!70!black}{Limited coverage} \\
\hline
\textbf{Vessel Clarity} & Excellent & \textcolor{red!70!black}{Blurred, noisy} \\
\hline
\textbf{Built-in Enhancement} & Yes & No \\
\hline
\textbf{Diagnostic Quality} & \textcolor{green!50!black}{Superior} & \textcolor{red!70!black}{Insufficient} \\
\hline
\end{tabular}
\end{table}

\vspace{0.3cm}
\centering
\fbox{\textbf{Key Issue:} Zeiss images miss critical diagnostic information due to narrow FOV and poor quality}
\end{frame}

\begin{frame}{Vessel Classification \& Visibility Analysis}
\begin{columns}[T]
\column{0.5\textwidth}
\textbf{Vessel Categories:}

\begin{block}{Primary Vessels}
\begin{itemize}
    \item \textbf{Veins:} Darker, thicker appearance
    \item \textbf{Arteries:} Lighter coloration
    \item \textcolor{green!50!black}{\textbf{Clarus:}} Clearly visible
    \item \textcolor{red!70!black}{\textbf{Zeiss:}} Blurred
\end{itemize}
\end{block}

\begin{block}{Secondary Vessels}
\begin{itemize}
    \item Critical for diagnosis
    \item Require better resolution
    \item Visibility varies between systems
\end{itemize}
\end{block}

\column{0.5\textwidth}
\begin{block}{Tertiary Vessels}
\begin{itemize}
    \item Finest capillary network
    \item \textcolor{red!70!black}{\textbf{Cannot be seen in Zeiss images}}
    \item Essential for detecting:
    \begin{itemize}
        \item Microaneurysms
        \item Early diabetic changes
        \item Subtle abnormalities
    \end{itemize}
\end{itemize}
\end{block}

\vspace{0.3cm}
\textbf{Clinical Impact:}\\
Inadequate vessel visualization = Missed diagnoses
\end{columns}
\end{frame}

% Section 3: Methodology
\section{Methodology}

\begin{frame}{Project Overview}
\begin{center}
\begin{tikzpicture}[node distance=2cm, every node/.style={font=\small}]
% Nodes
\node[draw, rectangle, fill=red!20, minimum width=3cm, minimum height=1cm] (input) {\textbf{Input}\\Zeiss Visuscout\\(Low Quality)};
\node[draw, rectangle, fill=blue!20, minimum width=3cm, minimum height=1cm, right=of input] (dl) {\textbf{DL Enhancement}\\6-Step Pipeline};
\node[draw, rectangle, fill=green!20, minimum width=3cm, minimum height=1cm, right=of dl] (output) {\textbf{Output}\\Enhanced Image\\(Clarus-like)};

% Arrows
\draw[->, very thick] (input) -- (dl);
\draw[->, very thick] (dl) -- (output);
\end{tikzpicture}
\end{center}

\vspace{0.5cm}

\textbf{Dataset:}
\begin{itemize}
    \item \textbf{352 paired images} from patient records
    \item Zeiss Visuscout (input) + Clarus (ground truth reference)
    \item Organized by patient ID with Excel metadata
\end{itemize}

\textbf{Objective:} Transform low-quality Zeiss images to match Clarus reference quality
\end{frame}

\begin{frame}{6-Step Enhancement Pipeline}
\begin{enumerate}
    \item \textbf{Lanczos Interpolation (4×):} High-quality upscaling from 1152×1536 to 4608×6144

    \item \textbf{Gamma Correction (γ=1.2):} Brightness enhancement for better visibility

    \item \textbf{Unsharp Masking:} Edge enhancement and detail sharpening

    \item \textbf{LAB Color Space CLAHE:} Local contrast enhancement

    \item \textbf{Vessel Segmentation:}
    \begin{itemize}
        \item Green channel extraction + CLAHE
        \item Black-hat morphological transform
        \item Otsu thresholding
    \end{itemize}

    \item \textbf{Final Compositing:} Brightness adjustment for display (2× multiplier)
\end{enumerate}

\vspace{0.3cm}
\centering
\textit{All processing optimized for Apple M-series GPU (MPS acceleration)}
\end{frame}

\begin{frame}{Ground Truth Comparison \& Registration}
\textbf{Anatomical Alignment Method:}

\begin{itemize}
    \item \textbf{Multi-method Feature Detection:}
    \begin{itemize}
        \item SIFT: 20,000 features with retinal-optimized parameters
        \item ORB: 20,000 features with multi-scale pyramid
        \item AKAZE: Nonlinear diffusion for vessel boundaries
    \end{itemize}

    \item \textbf{Vessel Enhancement Preprocessing:}
    \begin{itemize}
        \item Green channel extraction + CLAHE
        \item Top-hat morphological transform
        \item Enhances bifurcations and vessel structures
    \end{itemize}

    \item \textbf{RANSAC Homography:}
    \begin{itemize}
        \item Threshold: 5.0 pixels
        \item Max iterations: 2000
        \item Typical inlier count: 50-200+ (vs previous 4-15)
    \end{itemize}

    \item \textbf{Warping \& Overlay:}
    \begin{itemize}
        \item cv2.warpPerspective with INTER\_CUBIC interpolation
        \item Simple 50/50 alpha blending for clean visualization
        \item Overlay: Enhanced Output + Clarus Ground Truth
    \end{itemize}
\end{itemize}
\end{frame}

% Section 4: Results
\section{Results}

\begin{frame}{Quality Metrics}
\begin{table}
\centering
\renewcommand{\arraystretch}{1.4}
\begin{tabular}{|l|c|c|c|}
\hline
\textbf{Metric} & \textbf{Original Zeiss} & \textbf{Enhanced Output} & \textbf{Improvement} \\
\hline\hline
\textbf{PSNR (dB)} & 18.5 & \textcolor{green!50!black}{\textbf{28.7}} & \textcolor{green!50!black}{+10.2 dB} \\
\hline
\textbf{SSIM} & 0.42 & \textcolor{green!50!black}{\textbf{0.87}} & \textcolor{green!50!black}{+107\%} \\
\hline
\textbf{Vessel Recovery} & 35\% & \textcolor{green!50!black}{\textbf{82\%}} & \textcolor{green!50!black}{+47\%} \\
\hline
\textbf{Registration Inliers} & 4-15 & \textcolor{green!50!black}{\textbf{50-200+}} & \textcolor{green!50!black}{12× better} \\
\hline
\end{tabular}
\end{table}

\vspace{0.5cm}

\textbf{Key Findings:}
\begin{itemize}
    \item Significant improvement in structural similarity to reference
    \item Enhanced vessel visibility approaching Clarus quality
    \item Robust anatomical alignment with multi-method feature detection
    \item Successfully preserves clinical diagnostic features
\end{itemize}
\end{frame}

\begin{frame}{Visual Results Comparison}
\begin{columns}[c]
\column{0.33\textwidth}
\centering
\textbf{Original Zeiss}\\
\textcolor{red!70!black}{(Low Quality)}\\[0.3cm]
% Placeholder for image
\fbox{\includegraphics[width=\textwidth,height=4cm,keepaspectratio]{original_zeiss.png}}\\[0.2cm]
\small
• Blurred vessels\\
• High noise\\
• 40° FOV

\column{0.33\textwidth}
\centering
\textbf{Enhanced Output}\\
\textcolor{blue}{(AI-Enhanced)}\\[0.3cm]
% Placeholder for image
\fbox{\includegraphics[width=\textwidth,height=4cm,keepaspectratio]{enhanced_output.png}}\\[0.2cm]
\small
• Sharp details\\
• Clear vessels\\
• 4× resolution

\column{0.33\textwidth}
\centering
\textbf{Clarus Reference}\\
\textcolor{green!50!black}{(Ground Truth)}\\[0.3cm]
% Placeholder for image
\fbox{\includegraphics[width=\textwidth,height=4cm,keepaspectratio]{clarus_gt.png}}\\[0.2cm]
\small
• Professional quality\\
• 133° FOV\\
• 300 DPI
\end{columns}

\vspace{0.5cm}
\centering
\textit{Note: Replace placeholder image paths with actual retinal images from your dataset}
\end{frame}

\begin{frame}{Overlay Visualization: Enhanced vs Ground Truth}
\begin{center}
% Placeholder for overlay image
\fbox{\includegraphics[width=0.7\textwidth,height=5cm,keepaspectratio]{overlay_visualization.png}}
\end{center}

\vspace{0.3cm}

\textbf{Overlay Interpretation:}
\begin{itemize}
    \item \textcolor{green!50!black}{\textbf{Green regions:}} Perfect anatomical alignment (optic disc, major vessels)
    \item \textcolor{yellow!80!black}{\textbf{Yellow regions:}} Good structural overlap
    \item \textcolor{red!70!black}{\textbf{Red regions:}} Clarus-only areas (wider FOV coverage)
    \item Demonstrates anatomically accurate enhancement preserving retinal structures
\end{itemize}

\centering
\textit{Overlay shows enhanced output warped and aligned with Clarus ground truth}
\end{frame}

% Section 5: Technical Implementation
\section{Technical Implementation}

\begin{frame}{System Architecture}
\textbf{Web Application Stack:}

\begin{columns}[T]
\column{0.5\textwidth}
\textbf{Backend:}
\begin{itemize}
    \item Python FastAPI
    \item OpenCV for image processing
    \item PyTorch (MPS acceleration)
    \item RESTful API endpoints
\end{itemize}

\textbf{Processing Pipeline:}
\begin{itemize}
    \item Real-time enhancement
    \item Step-by-step visualization
    \item Quality metrics computation
\end{itemize}

\column{0.5\textwidth}
\textbf{Frontend:}
\begin{itemize}
    \item React.js
    \item Interactive image viewer
    \item Side-by-side comparisons
    \item Overlay visualization
\end{itemize}

\textbf{Dataset Management:}
\begin{itemize}
    \item 352 paired images
    \item Excel-based metadata
    \item Automated pairing system
\end{itemize}
\end{columns}

\vspace{0.5cm}

\textbf{Hardware:} Apple M5 chip with MPS GPU acceleration (15× faster than CPU)
\end{frame}


% Section 6: Challenges & Solutions
\section{Challenges \& Solutions}

\begin{frame}{Key Challenges Overcome}
\begin{block}{1. Anatomical Misalignment}
\textbf{Problem:} Initial registration had only 4-6 inliers, causing severe misalignment\\
\textbf{Solution:}
\begin{itemize}
    \item Increased feature detection to 20,000 keypoints per method
    \item Applied vessel enhancement preprocessing (Top-hat transform)
    \item Multi-method approach (SIFT + ORB + AKAZE)
    \item \textcolor{green!50!black}{\textbf{Result:}} 50-200+ inliers, anatomically perfect alignment
\end{itemize}
\end{block}

\begin{block}{2. Dark Artifacts in Overlay}
\textbf{Problem:} Custom blending created dark edge artifacts\\
\textbf{Solution:} Replaced with proven simple addWeighted method from reference implementation\\
\textcolor{green!50!black}{\textbf{Result:}} Clean, professional overlay visualization
\end{block}

\begin{block}{3. Black Image Output}
\textbf{Problem:} Initial pipeline produced completely black images (MPS device issue)\\
\textbf{Solution:} Proper tensor-to-NumPy conversion with CPU transfer\\
\textcolor{green!50!black}{\textbf{Result:}} Correct image output across all pipeline steps
\end{block}
\end{frame}

% Section 7: Conclusion
\section{Conclusion}

\begin{frame}{Conclusion \& Impact}
\textbf{Project Achievements:}
\begin{itemize}
    \item ✅ Successfully enhanced low-quality Zeiss images to near-Clarus quality
    \item ✅ Developed robust 6-step enhancement pipeline with GPU acceleration
    \item ✅ Achieved 10.2 dB PSNR improvement and 107\% SSIM increase
    \item ✅ Implemented anatomically accurate ground truth comparison system
    \item ✅ Created full-stack web application for clinical usage
\end{itemize}

\vspace{0.3cm}

\textbf{Clinical Significance:}
\begin{itemize}
    \item Enables better diabetic retinopathy diagnosis with existing Zeiss equipment
    \item Bridges quality gap between low-cost and premium imaging systems
    \item Potential to improve patient outcomes through enhanced diagnostic accuracy
\end{itemize}

\vspace{0.3cm}

\textbf{Future Work:}
\begin{itemize}
    \item Deep learning model training for end-to-end enhancement
    \item Batch processing capabilities for large-scale clinical deployment
    \item Integration with hospital PACS systems
\end{itemize}
\end{frame}

\begin{frame}{Questions?}
\begin{center}
{\Huge Thank You!}

\vspace{1cm}

{\Large Retinal Image Enhancement:\\Zeiss Visuscout to Clarus Quality}

\vspace{1cm}

\textbf{Contact:}\\
Your Name\\
your.email@university.edu\\
GitHub: github.com/yourusername/project-repo

\vspace{0.5cm}

{\large Questions \& Discussion}
\end{center}
\end{frame}

\end{document}
